\section{Method}
\label{sec:method}
\label{sec:formulation}

\subsection{Compilation Flow and Problem Definition}

Atomic positions at a layer boundary determine the execution conditions for current gates and the starting points of subsequent rearrangements. GA-LK optimizes transitions between consecutive two-qubit layers, jointly selecting gate sites and orientations, inter-layer residency, and return storage sites. Let $s_\ell$ denote the input state of target layer $\mathcal G_\ell$. It contains the mapping $\pi_\ell$, resident atoms in the entanglement zone, AOD resources, per-atom clocks, and gate dependencies. Two-qubit gates in $\mathcal G_\ell$ share no atoms. The corresponding rearrangements and gate execution produce the next layer's input state $s_{\ell+1}$.

GA-LK performs joint layer-boundary optimization during dynamic placement. Gate scheduling, initial layout, atom routing, and zoned architecture intermediate representation (ZAIR) code generation follow the zoned compilation flow~\cite{lin2025zac}. Scheduling partitions the input into alternating one- and two-qubit layers, and the initial layout assigns atoms to traps. Layer-boundary optimization compares the rearrangement batches and physical costs of candidate solutions. Code generation converts the selected solution into ZAIR instructions in execution order. Fig.~\ref{fig:overall-framework} shows these stages and their inter-layer state relationships.

\begin{figure*}[t]
  \centering
  \includegraphics[width=\textwidth]{../build/paper_zh/figures/overall_framework.pdf}
  \caption{End-to-end compilation flow for zoned neutral-atom circuits. (a) Input circuit and gate-execution zones; (b) layer scheduling; (c) initial mapping $\pi_0$. (d) Joint layer-boundary decisions produce candidate mapping $\pi'_\ell$. After the physical state transition, finite-horizon evaluation starts from $s_{\ell+1}$. (e) Compatible return and entry trajectories form rearrangement batches $\mathcal B_0$--$\mathcal B_2$. (f) ZAIR instructions are emitted in sequence and independently validated. $\mathcal E$ and $\mathcal S$ denote the entanglement and storage zones, respectively.}
  \label{fig:overall-framework}
\end{figure*}

Let $\mathcal Q$ be the logical atom set, and let $\mathcal S$ and $\mathcal E$ contain single-atom traps in the storage and entanglement zones, respectively. Define $\mathcal V=\mathcal S\cup\mathcal E$. The set $\mathcal P^{\rm gate}\subseteq\mathcal E\times\mathcal E$ contains paired gate sites in the entanglement zone. Injectivity of the boundary mapping permits at most one atom per trap:
\begin{equation}
\pi_\ell:\mathcal Q\rightarrow\mathcal V,\qquad
\pi_\ell(q)\ne\pi_\ell(q'),~q\ne q'.
\label{eq:mapping}
\end{equation}
For target gate $g_{\ell,j}=(u_j,v_j)\in\mathcal G_\ell$, the set $\Omega_{\ell,j}\subseteq\mathcal P^{\rm gate}$ specifies legal gate sites and orientations. The resulting candidate mapping $\pi'_{\ell}$ must also satisfy
\begin{equation}
\begin{aligned}
(\pi'_{\ell}(u_j),\pi'_{\ell}(v_j))&\in\Omega_{\ell,j},
&&\forall g_{\ell,j}\in\mathcal G_\ell,\\
|\pi_{\ell}^{\prime-1}(\mathcal E)|&\le C_{\mathcal E},
&|\pi_{\ell}^{\prime-1}(\mathcal S)|&\le C_{\mathcal S},
\end{aligned}
\label{eq:placement-constraints}
\end{equation}
where $C_{\mathcal E}$ and $C_{\mathcal S}$ are the zone capacities. The set $\mathcal D_\ell$ contains entanglement-zone atoms whose destinations must be decided at the current boundary. These include atoms absent from the target layer and those with no remaining gates. For a target layer containing one two-qubit gate, participating atoms already in the entanglement zone may also enter $\mathcal D_\ell$. Their candidate paths compare in-zone adjustment with temporary relocation through the storage zone. Raman single-qubit rotations update per-atom clocks according to gate dependencies~\cite{lin2025zac}.

\subsection{Joint Representation and Physical Feasible Domain}

Target gate sites and resident atoms share the entanglement zone, so both decisions must jointly satisfy occupancy constraints. A layer-boundary solution encodes gate-site and residency choices as
\begin{equation}
\mathbf x_\ell=(z_{\ell,1},\ldots,z_{\ell,m},b_{\ell,1},\ldots,b_{\ell,k}),
\label{eq:chromosome}
\end{equation}
where $m=|\mathcal G_\ell|$ and $k=|\mathcal D_\ell|$. Each $z_{\ell,j}\in\{0,\ldots,|\Omega_{\ell,j}|-1\}$ selects a gate site and orientation; $b_{\ell,i}=0$ denotes residency, whereas $b_{\ell,i}=1$ denotes return to storage. The two encoding segments determine the target gate sites and the set of returning atoms. A storage-site assignment completes the candidate's resulting state. Fig.~\ref{fig:joint-ga} relates the encoding to trajectories using two target gates and the return of $q_3$.

\begin{figure}[!t]
  \centering
  \resizebox{\columnwidth}{!}{% Shared vector-figure vocabulary for the Chinese IEEE manuscript.
% Every figure in this directory inputs this file, so the manuscript only
% needs the standard TikZ package and the libraries already loaded by
% paper_zh.tex.  Keep the guard: figures may be input more than once while
% editing or when a stand-alone smoke-test document is built.
\ifdefined\GALKFigureStyleLoaded\else
\def\GALKFigureStyleLoaded{1}

% Baseline-derived visual tokens: ICCAD uses a restrained blue/orange device
% palette, while ZAC uses conventional black axes and a small fixed plot set.
\definecolor{galkInk}{HTML}{231F20}
\definecolor{galkMuted}{HTML}{666666}
\definecolor{galkGrid}{HTML}{D6D6D6}
\definecolor{galkPaper}{HTML}{FFFFFF}
% Semantic colors are fixed across every figure: orange=entanglement zone,
% blue=storage zone, black=atom identity, red=invalidity, and blue dashes=future
% prediction.  Plot and search-group colours are deliberately independent of
% the two physical-zone colours.
\definecolor{galkStorage}{HTML}{0064BC}
\definecolor{galkStorageFill}{HTML}{E8F1F7}
\definecolor{galkEntangle}{HTML}{E37323}
\definecolor{galkEntangleFill}{HTML}{F9ECE3}
\definecolor{galkResident}{HTML}{485CB3}
\definecolor{galkResidentFill}{HTML}{EEF0F8}
\definecolor{galkGood}{HTML}{A4AF07}
\definecolor{galkBad}{HTML}{C84E45}
\definecolor{galkLook}{HTML}{0064BC}
\definecolor{galkData}{HTML}{315E8A}
\definecolor{galkTail}{HTML}{B44B27}
\definecolor{galkGroupOne}{HTML}{4C78A8}
\definecolor{galkGroupTwo}{HTML}{7A7A7A}
\definecolor{galkGroupThree}{HTML}{9C755F}

\tikzset{
  % Shared engineering-figure contract.  Physical snapshots must place an
  % orange entanglement band above a blue storage field; only their scale may
  % change.  Labels are aligned at the upper-left of both regions.
  galk/base/.style={
    font=\rmfamily\fontsize{9.0pt}{10.0pt}\selectfont,
    text=galkInk,
    line cap=butt,
    line join=miter,
    >=Stealth
  },
  galk/panel title/.style={
    font=\rmfamily\bfseries\fontsize{9.2pt}{10.2pt}\selectfont,
    text=galkInk,
    anchor=west
  },
  galk/label/.style={
    font=\rmfamily\fontsize{9.0pt}{10.0pt}\selectfont,
    text=galkInk
  },
  galk/small label/.style={
    font=\rmfamily\fontsize{8.7pt}{9.7pt}\selectfont,
    text=galkInk
  },
  galk/zone label/.style={
    font=\rmfamily\bfseries\fontsize{8.7pt}{9.7pt}\selectfont,
    text=galkInk,
    anchor=north west,
    inner sep=.5pt
  },
  galk/entanglement zone/.style={
    draw=galkEntangle,
    fill=galkEntangleFill,
    line width=.40pt
  },
  galk/storage zone/.style={
    draw=galkStorage,
    fill=galkStorageFill,
    line width=.40pt
  },
  galk/site/.style={
    circle,
    draw=galkInk!68,
    fill=white,
    line width=.36pt,
    minimum size=1.35mm,
    inner sep=0pt
  },
  galk/occupied atom/.style={
    circle,
    draw=galkInk,
    fill=galkInk,
    line width=.36pt,
    minimum size=1.55mm,
    inner sep=0pt
  },
  galk/target site/.style={
    circle,
    draw=galkInk!60,
    fill=galkInk!10,
    line width=.34pt,
    minimum size=1.55mm,
    inner sep=0pt
  },
  galk/q label/.style={
    font=\rmfamily\itshape\fontsize{8.5pt}{9.5pt}\selectfont,
    text=galkInk,
    inner sep=.25pt
  },
  galk/rydberg pair/.style={
    draw=galkInk!62,
    densely dashed,
    line width=.36pt
  },
  galk/accepted path/.style={
    -{Stealth[length=.95mm,width=.62mm]},
    draw=galkInk,
    line width=.48pt
  },
  galk/alternative path/.style={
    -{Stealth[length=.95mm,width=.62mm]},
    draw=galkMuted,
    densely dashed,
    line width=.48pt
  },
  galk/future path/.style={
    -{Stealth[length=.95mm,width=.62mm]},
    draw=galkLook,
    densely dashed,
    line width=.48pt
  },
  galk/invalid path/.style={
    -{Stealth[length=.95mm,width=.62mm]},
    draw=galkBad,
    densely dashed,
    line width=.50pt
  }
}
\fi

\begin{tikzpicture}[
  galk/base,x=1.04cm,y=1cm,
  panel title/.style={galk/panel title,font=\rmfamily\bfseries\fontsize{9.6pt}{10.6pt}\selectfont},
  label/.style={font=\rmfamily\fontsize{9.2pt}{10.2pt}\selectfont},
  note/.style={font=\rmfamily\itshape\fontsize{9.0pt}{10.0pt}\selectfont,
               text=galkMuted},
  tiny/.style={font=\rmfamily\fontsize{9.0pt}{10.0pt}\selectfont,
               text=galkInk},
  rule/.style={draw=galkInk!30,line width=.38pt},
  flow/.style={-{Stealth[length=1.0mm,width=.66mm]},draw=galkInk,line width=.50pt},
  cell/.style={draw=galkInk!72,fill=white,line width=.38pt,
               minimum width=.77cm,minimum height=.43cm,inner sep=.5pt},
  gate cell/.style={cell,draw=galkEntangle,fill=white},
  bit cell/.style={cell,draw=galkResident,fill=galkResidentFill},
  trap/.style={galk/site},
  atom/.style={galk/occupied atom,minimum size=1.85mm},
  target/.style={galk/target site,minimum size=1.85mm},
  q label/.style={galk/q label,font=\rmfamily\itshape\fontsize{9.0pt}{10.0pt}\selectfont}
]
  \path[use as bounding box] (0,.40) rectangle (8.35,7.35);
  \draw[rule] (.10,5.12)--(8.25,5.12);
  \draw[rule] (.10,4.05)--(8.25,4.05);

  % ---------------------------------------------------------------------
  % A joint chromosome and the induced injective storage assignment.
  \node[panel title] at (.12,7.15) {(a) Encoding a placement candidate};
  \node[note] at (1.93,6.83) {chromosome $\mathbf x_\ell$};
  \foreach \i/\head/\value/\cellstyle in {
      0/{$z_{\ell,1}$}/{$T_1^{+}$}/gate cell,
      1/{$z_{\ell,2}$}/{$T_2^{+}$}/gate cell,
      2/{$b_{\ell,1}$}/{1}/bit cell} {
    \node[tiny] at (.76+1.17*\i,6.46) {\head};
    \node[\cellstyle] at (.76+1.17*\i,6.10) {\value};
  }
  \node[tiny] at (.76,5.68) {$(q_0,q_2)$};
  \node[tiny] at (1.93,5.68) {$(q_1,q_4)$};
  \node[tiny] at (3.10,5.68) {$q_3$};
  \node[note] at (1.93,5.30) {gate atoms: left to right};

  \draw[flow] (3.58,6.10)--(4.12,6.10);
  \node[label] at (6.14,6.90) {$K$-best storage assignment};
  \node[tiny,anchor=west] at (4.25,6.53)
    {$q_3$ to storage};
  \node[atom] at (4.47,6.12) {};
  \node[q label,anchor=north] at (4.47,5.99) {$q_3$};
  \foreach \y/\site in {6.55/{$s_1$},6.10/{$s_2$},5.65/{$s_3$}} {
    \node[trap,minimum size=2.05mm] at (7.39,\y) {};
    \node[tiny,anchor=west] at (7.56,\y) {\site};
  }
  \draw[galkInk!34,line width=.34pt] (4.57,6.12)--(7.29,6.55);
  \draw[galkInk!34,line width=.34pt] (4.57,6.12)--(7.29,6.10);
  \draw[galkInk!34,line width=.34pt] (4.57,6.12)--(7.29,5.65);
  \draw[galkInk,line width=.62pt] (4.57,6.12)--(7.29,6.55);
  \node[label,anchor=east] at (8.15,5.30)
    {joint candidate $(\mathbf x_\ell,a_\ell)$};

  % ---------------------------------------------------------------------
  % Enumeration and genetic search share the same executable evaluation.
  \node[panel title] at (.12,4.88) {(b) Search with a common objective};
  \node[tiny] at (2.20,4.39) {Enumeration / Genetic search};
  \draw[flow] (4.46,4.39)--(4.93,4.39);
  \node[tiny] at (6.67,4.39) {Candidate evaluation};

  % ---------------------------------------------------------------------
  % Release the target sites, including temporary storage for the q1/q2 swap.
  \def\routeInitial{0/.36/1.91,1/.58/1.91,2/1.02/1.91,3/1.24/1.91,4/.63/.70}
  \def\routeStaged{0/.36/1.91,1/.28/1.08,2/1.02/1.91,3/1.33/1.08,4/.63/.70}
  \def\routeFinal{0/.36/1.91,1/1.02/1.91,2/.58/1.91,3/1.33/1.08,4/1.24/1.91}
  \def\routeMoves{0/1/.58/1.91/.28/1.08,0/3/1.24/1.91/1.33/1.08,1/2/1.02/1.91/.58/1.91,2/1/.28/1.08/1.02/1.91,3/4/.63/.70/1.24/1.91}
  \begin{scope}[shift={(0,.30)},x=1.18cm]
  \node[panel title] at (.12,3.40) {(c) Decoding and ordered execution};

  \foreach \left in {.18,2.93,5.68} {
    \path[galk/entanglement zone]
      (\left,1.63) rectangle (\left+1.60,2.80);
    \path[galk/storage zone]
      (\left,.20) rectangle (\left+1.60,1.55);
    \node[tiny] at (\left+.47,2.55) {$T_1$};
    \node[tiny] at (\left+1.13,2.55) {$T_2$};
    \draw[galk/rydberg pair] (\left+.47,1.91) ellipse (.20 and .13);
    \draw[galk/rydberg pair] (\left+1.13,1.91) ellipse (.20 and .13);
    \foreach \dx in {.36,.58,1.02,1.24}
      \node[trap] at (\left+\dx,1.91) {};
    \foreach \dx in {.28,.63,.98,1.33}
      \foreach \yy in {.70,1.08}
        \node[trap] at (\left+\dx,\yy) {};
  }
  \node[label] at (.98,3.00) {$\mathbf{s}_\ell$};
  \node[label] at (3.73,3.00) {After staging};
  \node[label] at (6.48,3.00) {$\mathbf{s}_{\ell+1}$};

  % Source state before G_l placement: two occupied pairs and q4 in storage.
  \foreach \qq/\xx/\yy in \routeInitial
    \node[atom] at (.18+\xx,\yy) {};
  \node[q label,font=\fontsize{8.5pt}{9.5pt}\selectfont] at (.65,2.20) {$q_0,q_1$};
  \node[q label,font=\fontsize{8.5pt}{9.5pt}\selectfont] at (1.31,2.20) {$q_2,q_3$};
  \draw[galkInk,line width=.46pt] (.54,1.91)--(.76,1.91);
  \draw[galkInk,line width=.46pt] (1.20,1.91)--(1.42,1.91);
  \node[q label,anchor=north] at (.81,.58) {$q_4$};

  % Intermediate state: q1 is temporarily stored; q3 remains in storage.
  \foreach \qq/\xx/\yy in \routeStaged
    \node[atom] at (2.93+\xx,\yy) {};
  \node[q label] at (3.40,2.20) {$q_0$};
  \node[q label] at (4.06,2.20) {$q_2$};
  \node[q label,anchor=north] at (3.56,.58) {$q_4$};
  \node[q label,anchor=south] at (3.21,1.20) {$q_1$};
  \node[q label,anchor=south east] at (4.44,1.20) {$q_3@s_1$};

  % Target state: q0-q2 and q1-q4 occupy the selected gate pairs.
  \foreach \qq/\xx/\yy in \routeFinal
    \node[atom] at (5.68+\xx,\yy) {};
  \node[q label,font=\fontsize{8.5pt}{9.5pt}\selectfont] at (6.15,2.20) {$q_0,q_2$};
  \node[q label,font=\fontsize{8.5pt}{9.5pt}\selectfont] at (6.81,2.20) {$q_1,q_4$};
  \draw[galkInk,line width=.46pt] (6.04,1.91)--(6.26,1.91);
  \draw[galkInk,line width=.46pt] (6.70,1.91)--(6.92,1.91);
  % The same s_1 trap (storage row 1, column 4) remains occupied by q3.
  \node[q label,anchor=south east] at (7.19,1.20) {$q_3@s_1$};

  \draw[flow] (1.91,1.58)--(2.80,1.58);
  \node[tiny,align=center] at (2.35,2.20) {Return to\\storage};
  \node[tiny] at (2.35,1.30) {$B_0$};
  \draw[flow] (4.66,1.58)--(5.55,1.58);
  \node[tiny,align=center] at (5.10,2.20) {Place gate\\atoms};
  \node[tiny] at (5.10,1.30) {$B_1$--$B_3$};
  \end{scope}
\end{tikzpicture}
}
  \caption{Joint placement at a two-qubit layer boundary. (a) $z_{\ell,1}$ and $z_{\ell,2}$ select gate sites and orientations. Here, $\mathcal D_\ell=\{q_3\}$, and $b_{\ell,1}=1$ includes $q_3$ in the return set; injective assignment $a_\ell(q_3)=s_1$ selects its storage site. (b) Exact enumeration and genetic search share the same encoding space. (c) The candidate first returns $q_3$ to $s_1$, updates occupancy, and then places target-layer participants at the selected gate sites. Each phase checks AOD row--column relationships, occupancy ordering, and unintended intersections. $\mathcal E$ and $\mathcal S$ denote the entanglement and storage zones, respectively.}
  \label{fig:joint-ga}
\end{figure}

Genetic search determines gate-site and residency combinations, while a matching subproblem assigns distinct storage sites to returning atoms. For each atom $q$ in the return set $\mathcal R(\mathbf x_\ell)$, candidate storage sites come from three neighborhoods. These are defined by its current entanglement site, initial storage site, and visible future partners, giving the bounded candidate domain
\begin{equation}
\mathcal S_q^{\rm cand}=\operatorname{Top}
\left(\mathcal S_q^{\rm near}\cup
\mathcal S_q^{\rm home}\cup
\mathcal S_q^{\rm future}\right),
\label{eq:storage-domain}
\end{equation}
where the three subsets correspond to these neighborhoods in order. The operator $\operatorname{Top}(\cdot)$ retains a limited number of sites according to a deterministic proxy cost, bounding the matching problem size. Linear assignment subproblems and Murty branching yield the first $K$ injective assignments~\cite{murty1968assignment}:
\begin{equation}
\begin{aligned}
\mathcal A_K(\mathbf x_\ell)=K\text{-Best}\{a:\mathcal R(\mathbf x_\ell)\!\rightarrow\!\mathcal S\mid{}&
a(q)\in\mathcal S_q^{\rm cand},\\[-1mm]
&a(q)\ne a(q'),~\forall q\ne q'\}.
\end{aligned}
\label{eq:kbest}
\end{equation}
Each assignment $a\in\mathcal A_K$ combines with $\mathbf x_\ell$ to form a complete candidate. The proxy cost serves only to generate candidates. Final ordering uses the physical state-transition cost and~\eqref{eq:lexicographic}.

Physical constraints fall into four groups: occupancy injectivity, zone capacity, and source-release and destination-occupancy ordering; gate sites and dependencies; AOD capacity, row--column ordering, and preservation; and unintended intersections. Their indicator functions are $\chi_{\rm occ}$, $\chi_{\rm gate}$, $\chi_{\rm AOD}$, and $\chi_{\rm int}$, respectively. The physical feasible domain is
\begin{equation}
\mathcal X_\ell^{\rm phys}=\{(\mathbf x,a)\mid
a\in\mathcal A_K(\mathbf x),
\chi_{\rm occ}\chi_{\rm gate}\chi_{\rm AOD}\chi_{\rm int}=1\}.
\label{eq:physical-domain}
\end{equation}
If a nonparticipating atom occupies a selected gate site, the candidate must also determine its relocation. Target-layer participants may temporarily use available storage sites to clear entry paths. Candidate decoding determines these auxiliary position adjustments; their trap transfers, transport, and clock increments all contribute to the physical cost. A chromosome is discarded if all its storage-site assignments are infeasible.

\subsection{Physical Cost and Parallel Rearrangement}

Physical cost depends on rearrangement order and the grouping of parallel trajectories. To clear target gate sites, each candidate first removes returning atoms and gate-site blockers from the entanglement zone. It then moves target-layer participants to the selected gate sites. Each phase forms its own parallel rearrangement batches, with occupancy updated after completed movements. Positions, resource occupancy, and per-atom clocks follow the dependencies of rearrangement and gate execution. For a feasible candidate $(\mathbf x,a)$, the state transition is
\begin{equation}
(s_{\ell+1},\Gamma_\ell)=
\mathcal T_\ell(s_\ell;\mathbf x,a),
\label{eq:state-transition}
\end{equation}
where $\Gamma_\ell=(\mathcal B_{\ell,1},\ldots,\mathcal B_{\ell,B_\ell})$ lists rearrangement batches in execution order. The state $s_{\ell+1}$ follows the target layer's Rydberg pulse and corresponding single-qubit gates. This transition includes idle-atom excitation for nonparticipating residents and coherence loss for each atom.

Parallel grouping uses a conflict graph $\mathcal H_\phi=(\mathcal V_\phi,\mathcal E_\phi^{\rm AOD}\cup\mathcal E_\phi^{\rm int})$, whose vertices represent atom trajectories. Edges in $\mathcal E_\phi^{\rm AOD}$ connect trajectories that violate AOD ordering. Edges in $\mathcal E_\phi^{\rm int}$ connect trajectories whose joint activation creates an unintended-intersection conflict. A candidate is infeasible if an individual trajectory violates static occupancy or intersection constraints. DSATUR coloring supplies initial groups for the remaining conflicts~\cite{brelaz1979coloring}. Trajectories of one color may share a batch, while different colors execute separately. Each group is then checked against source-release and destination-occupancy ordering and split when necessary. Let $\tau(\nu)$ and $d(\nu)$ denote the execution time and transport distance of trajectory $\nu$, respectively. Then
\begin{equation}
T_\ell^{\rm rr}=\sum_{b=1}^{B_\ell}
\max_{\nu\in\mathcal B_{\ell,b}}\tau(\nu),\qquad
D_\ell^{\rm rr}=\sum_{b=1}^{B_\ell}
\sum_{\nu\in\mathcal B_{\ell,b}}d(\nu).
\label{eq:rearrangement-metrics}
\end{equation}
The batch count measures the number of sequential rearrangement phases. $T_\ell^{\rm rr}$ and $D_\ell^{\rm rr}$ measure rearrangement time and total transport distance, respectively.

For a given circuit, single- and two-qubit gate errors are identical across candidates. Boundary optimization compares each candidate's increment in physical negative log fidelity relative to the preceding state:
\begin{equation}
\begin{aligned}
C_\ell^{\rm cur}={}&-\Delta N_{\rm tran}\log f_{\rm tran}
-\Delta N_{\rm exc}\log f_{\rm exc}\\
&+\sum_{q\in\mathcal Q}\log\frac{1-t_q^{-}/T_2}{1-t_q^{+}/T_2}.
\end{aligned}
\label{eq:current-cost}
\end{equation}
Here, $t_q^{-}$ and $t_q^{+}$ denote the accumulated idle times before and after the state transition, respectively.

\subsection{Decayed Multilayer Look-Ahead}

Decayed multilayer physical look-ahead includes interactions several layers ahead in the residency evaluation. It estimates future physical costs from each candidate's resulting state and assigns higher weights to nearer layers. For maximum horizon $H_{\max}$, the weight at layer distance $d$ is
\begin{equation}
w_d=\alpha\rho^{d-1},\qquad 1\le d\le H_{\max}.
\label{eq:decay}
\end{equation}
Let $H_{\rm cap}\le H_{\max}$ denote the resource-dependent horizon cap. The finite horizon contains future two-qubit layers that satisfy the layer-distance, remaining-layer, resource-cap, and decay-threshold conditions. Its set of layer distances is
\begin{equation}
\mathcal D_\ell^{\rm vis}=\{d\mid 1\le d\le H_{\rm cap},~
\mathcal G_{\ell+d}\text{ exists},~\rho^{d-1}\ge\epsilon\}.
\label{eq:visible-depths}
\end{equation}
The main configuration uses $(H_{\max},\alpha,\rho,\epsilon)=(8,0.5,0.7,0.05)$, with $H_{\rm cap}$ reduced as the problem size increases. Fig.~\ref{fig:physical-lookahead} illustrates how the candidate's resulting state affects next-layer transport and how weights vary with layer distance.

\begin{figure}[!t]
  \centering
  \resizebox{\columnwidth}{!}{% Shared vector-figure vocabulary for the Chinese IEEE manuscript.
% Every figure in this directory inputs this file, so the manuscript only
% needs the standard TikZ package and the libraries already loaded by
% paper_zh.tex.  Keep the guard: figures may be input more than once while
% editing or when a stand-alone smoke-test document is built.
\ifdefined\GALKFigureStyleLoaded\else
\def\GALKFigureStyleLoaded{1}

% Baseline-derived visual tokens: ICCAD uses a restrained blue/orange device
% palette, while ZAC uses conventional black axes and a small fixed plot set.
\definecolor{galkInk}{HTML}{231F20}
\definecolor{galkMuted}{HTML}{666666}
\definecolor{galkGrid}{HTML}{D6D6D6}
\definecolor{galkPaper}{HTML}{FFFFFF}
% Semantic colors are fixed across every figure: orange=entanglement zone,
% blue=storage zone, black=atom identity, red=invalidity, and blue dashes=future
% prediction.  Plot and search-group colours are deliberately independent of
% the two physical-zone colours.
\definecolor{galkStorage}{HTML}{0064BC}
\definecolor{galkStorageFill}{HTML}{E8F1F7}
\definecolor{galkEntangle}{HTML}{E37323}
\definecolor{galkEntangleFill}{HTML}{F9ECE3}
\definecolor{galkResident}{HTML}{485CB3}
\definecolor{galkResidentFill}{HTML}{EEF0F8}
\definecolor{galkGood}{HTML}{A4AF07}
\definecolor{galkBad}{HTML}{C84E45}
\definecolor{galkLook}{HTML}{0064BC}
\definecolor{galkData}{HTML}{315E8A}
\definecolor{galkTail}{HTML}{B44B27}
\definecolor{galkGroupOne}{HTML}{4C78A8}
\definecolor{galkGroupTwo}{HTML}{7A7A7A}
\definecolor{galkGroupThree}{HTML}{9C755F}

\tikzset{
  % Shared engineering-figure contract.  Physical snapshots must place an
  % orange entanglement band above a blue storage field; only their scale may
  % change.  Labels are aligned at the upper-left of both regions.
  galk/base/.style={
    font=\rmfamily\fontsize{9.0pt}{10.0pt}\selectfont,
    text=galkInk,
    line cap=butt,
    line join=miter,
    >=Stealth
  },
  galk/panel title/.style={
    font=\rmfamily\bfseries\fontsize{9.2pt}{10.2pt}\selectfont,
    text=galkInk,
    anchor=west
  },
  galk/label/.style={
    font=\rmfamily\fontsize{9.0pt}{10.0pt}\selectfont,
    text=galkInk
  },
  galk/small label/.style={
    font=\rmfamily\fontsize{8.7pt}{9.7pt}\selectfont,
    text=galkInk
  },
  galk/zone label/.style={
    font=\rmfamily\bfseries\fontsize{8.7pt}{9.7pt}\selectfont,
    text=galkInk,
    anchor=north west,
    inner sep=.5pt
  },
  galk/entanglement zone/.style={
    draw=galkEntangle,
    fill=galkEntangleFill,
    line width=.40pt
  },
  galk/storage zone/.style={
    draw=galkStorage,
    fill=galkStorageFill,
    line width=.40pt
  },
  galk/site/.style={
    circle,
    draw=galkInk!68,
    fill=white,
    line width=.36pt,
    minimum size=1.35mm,
    inner sep=0pt
  },
  galk/occupied atom/.style={
    circle,
    draw=galkInk,
    fill=galkInk,
    line width=.36pt,
    minimum size=1.55mm,
    inner sep=0pt
  },
  galk/target site/.style={
    circle,
    draw=galkInk!60,
    fill=galkInk!10,
    line width=.34pt,
    minimum size=1.55mm,
    inner sep=0pt
  },
  galk/q label/.style={
    font=\rmfamily\itshape\fontsize{8.5pt}{9.5pt}\selectfont,
    text=galkInk,
    inner sep=.25pt
  },
  galk/rydberg pair/.style={
    draw=galkInk!62,
    densely dashed,
    line width=.36pt
  },
  galk/accepted path/.style={
    -{Stealth[length=.95mm,width=.62mm]},
    draw=galkInk,
    line width=.48pt
  },
  galk/alternative path/.style={
    -{Stealth[length=.95mm,width=.62mm]},
    draw=galkMuted,
    densely dashed,
    line width=.48pt
  },
  galk/future path/.style={
    -{Stealth[length=.95mm,width=.62mm]},
    draw=galkLook,
    densely dashed,
    line width=.48pt
  },
  galk/invalid path/.style={
    -{Stealth[length=.95mm,width=.62mm]},
    draw=galkBad,
    densely dashed,
    line width=.50pt
  }
}
\fi

\begin{tikzpicture}[galk/base,x=1cm,y=1cm,
  panel title/.style={galk/panel title},
  small label/.style={galk/small label},
  zone label/.style={galk/zone label},
  site/.style={galk/site},
  occupied/.style={galk/occupied atom},
  target/.style={galk/target site},
  route/.style={galk/accepted path},
  graph vertex/.style={circle,draw=galkInk,line width=.42pt,
                       minimum size=2.25mm,inner sep=0pt}
]
  \path[use as bounding box] (0,-.22) rectangle (8.50,8.22);

  % (a) Physical trajectories, their conflict graph, and the resulting groups.
  \node[panel title] at (.08,7.98) {(a) Compatibility at $\pi_\ell\!\to\!\pi_{\ell+1}$};
  \node[small label] at (1.50,7.62) {trajectories};
  \node[small label,anchor=east] at (.53,7.50) {$\pi_\ell$};
  \node[small label,anchor=west] at (2.54,7.50) {$\pi_{\ell+1}$};

  \coordinate (s1) at (.58,7.20);
  \coordinate (s2) at (.58,6.82);
  \coordinate (s3) at (.58,6.42);
  \coordinate (s4) at (.58,6.04);
  \coordinate (t1) at (2.48,6.82);
  \coordinate (t2) at (2.48,7.20);
  \coordinate (t3) at (2.48,6.04);
  \coordinate (t4) at (2.48,6.42);
  \foreach \i in {1,2,3,4} {
    \node[occupied] at (s\i) {};
    \node[target] at (t\i) {};
    \node[small label,anchor=east] at ([xshift=-1.4pt]s\i) {$q_{\i}$};
    \node[small label,anchor=west] at ([xshift=1.4pt]t\i) {$q_{\i}$};
  }
  \draw[route] (s1)--(t1);
  \draw[route] (s2)--(t2);
  \draw[route] (s3)--(t3);
  \draw[route] (s4)--(t4);

  \draw[-{Stealth[length=.85mm,width=.55mm]},galkInk,line width=.42pt]
    (2.95,6.61)--(3.35,6.61);
  \node[small label] at (4.25,7.62) {conflict graph};

  \node[graph vertex,fill=galkGroupTwo!18] (v1) at (3.68,7.18) {};
  \node[graph vertex,fill=galkGroupOne!18] (v2) at (4.15,6.43) {};
  \node[graph vertex,fill=galkGroupTwo!18] (v3) at (4.83,6.94) {};
  \node[graph vertex,fill=galkGroupOne!18] (v4) at (4.89,6.10) {};
  \node[small label,anchor=east] at ([xshift=-1.5pt]v1.west) {$q_1$};
  \node[small label,anchor=east] at ([xshift=-1.5pt]v2.west) {$q_2$};
  \node[small label,anchor=west] at ([xshift=1.5pt]v3.east) {$q_3$};
  \node[small label,anchor=west] at ([xshift=1.5pt]v4.east) {$q_4$};
  \draw[galkBad,line width=.46pt] (v1)--(v2);
  \draw[galkBad,line width=.46pt] (v3)--(v4);
  \node[small label] at (4.25,5.73)
    {$\{q_1,q_2\},\ \{q_3,q_4\}$};

  \draw[-{Stealth[length=.85mm,width=.55mm]},galkInk,line width=.42pt]
    (5.34,6.61)--(5.77,6.61);
  \node[small label] at (6.88,7.62) {DSATUR coloring};
  \node[small label,anchor=west] at (5.92,7.20)
    {$\mathcal B_1=\{q_2,q_4\}$};
  \node[small label,anchor=west] at (5.92,6.67)
    {$\mathcal B_2=\{q_1,q_3\}$};
  \draw[galkGroupOne,line width=1.05pt] (5.92,7.03)--(6.32,7.03);
  \draw[galkGroupTwo,line width=1.05pt] (5.92,6.50)--(6.32,6.50);
  \node[small label,anchor=west,align=left] at (5.92,6.10)
    {$q_0$: stationary\\no trajectory};

  \draw[galkInk!32,line width=.35pt] (.08,5.47)--(8.20,5.47);

  % (b) Equal immediate costs may lead to different finite-horizon costs.
  \node[panel title] at (.08,5.31) {(b) Dependence on the post-layer state};
  \node[small label] at (4.17,4.96)
    {$C_{\ell}^{\rm cur}(A)=C_{\ell}^{\rm cur}(B)$};

  \foreach \xzero in {.42,4.65} {
    \path[galk/entanglement zone]
      (\xzero,3.70) rectangle (\xzero+3.22,4.64);
    \path[galk/storage zone]
      (\xzero,2.08) rectangle (\xzero+3.22,3.54);
    \node[zone label,fill=galkEntangleFill]
      at (\xzero+.08,4.58) {$\mathcal E$};
    \node[zone label,fill=galkStorageFill]
      at (\xzero+.08,3.48) {$\mathcal S$};
    \foreach \dx in {.38,.98,1.58,2.18,2.78} {
      \foreach \yy in {2.38,2.80,3.22} {
        \node[site] at (\xzero+\dx,\yy) {};
      }
    }
    \foreach \px in {.72,2.18} {
      \node[site] at (\xzero+\px,3.95) {};
      \node[site] at (\xzero+\px+.26,3.95) {};
      \draw[galk/rydberg pair]
        (\xzero+\px+.13,3.95) ellipse (.21 and .13);
    }
  }
  \node[small label] at (1.10,4.94)
    {$A:\ s_{\ell+1}$};
  \node[small label] at (7.18,4.94)
    {$B:\ s_{\ell+1}$};

  % Candidate A: q and r occupy the current gate pair.  The next layer moves
  % q and p to the second pair after r returns to storage.
  \node[occupied] (ar) at (1.14,3.95) {};
  \node[occupied] (aq) at (1.40,3.95) {};
  \node[target] (aqt) at (2.60,3.95) {};
  \node[target] (apt) at (2.86,3.95) {};
  \node[occupied] (ap) at (3.20,2.38) {};
  \node[target] (art) at (2.00,3.22) {};
  \node[galk/q label,anchor=north east] at ([xshift=-2.8pt,yshift=-1.4pt]ar) {$r$};
  \node[galk/q label,anchor=north west] at ([xshift=2.8pt,yshift=-1.4pt]aq) {$q$};
  \node[galk/q label,anchor=west] at ([xshift=3.2pt]ap) {$p$};
  \node[small label] at (1.27,4.18) {$\mathrm{CZ}(q,r)$};
  \node[small label] at (2.73,4.18) {$\mathrm{CZ}(q,p)$};
  \draw[galk/alternative path]
    (ar) .. controls (1.45,3.72) and (1.76,3.48) .. (art);
  \draw[galk/future path] (aq)--(aqt);
  \draw[galk/future path]
    (ap) .. controls (3.25,3.12) and (3.10,3.57) ..
    node[pos=.44,left,small label] {$d_A$} (apt);

  % Candidate B has the same current gate pair and target pair, but places p
  % nearer the target used by CZ(q,p).
  \node[occupied] (br) at (5.37,3.95) {};
  \node[occupied] (bq) at (5.63,3.95) {};
  \node[target] (bqt) at (6.83,3.95) {};
  \node[target] (bpt) at (7.09,3.95) {};
  \node[occupied] (bp) at (6.83,3.22) {};
  \node[target] (brt) at (6.23,3.22) {};
  \node[galk/q label,anchor=north east] at ([xshift=-2.8pt,yshift=-1.4pt]br) {$r$};
  \node[galk/q label,anchor=north west] at ([xshift=2.8pt,yshift=-1.4pt]bq) {$q$};
  \node[galk/q label,anchor=west] at ([xshift=3.2pt]bp) {$p$};
  \node[small label] at (5.50,4.18) {$\mathrm{CZ}(q,r)$};
  \node[small label] at (6.96,4.18) {$\mathrm{CZ}(q,p)$};
  \draw[galk/alternative path]
    (br) .. controls (5.68,3.72) and (5.99,3.48) .. (brt);
  \draw[galk/future path] (bq)--(bqt);
  \draw[galk/future path]
    (bp)--node[pos=.70,right,small label,fill=white,inner sep=.2pt] {$d_B$} (bpt);

  \node[small label] at (4.17,1.82)
    {$d_B<d_A\ \Rightarrow\ C_{\ell,1}^{\rm gate}(B)<C_{\ell,1}^{\rm gate}(A)$};

  \draw[galkInk!32,line width=.35pt] (.08,1.55)--(8.20,1.55);

  % (c) Conventional axis representation of the geometric decay.
  \node[panel title] at (.08,1.31) {(c) Decayed layer weights};
  \draw[-{Stealth[length=.8mm,width=.52mm]},galkInk,line width=.42pt]
    (.72,.27)--(5.14,.27) node[right,small label] {$d$};
  \draw[-{Stealth[length=.8mm,width=.52mm]},galkInk,line width=.42pt]
    (.72,.27)--(.72,1.05);
  \node[small label,rotate=90] at (.49,.69) {$w_d$};
  \draw[galkData,line width=.55pt]
    (1.23,1.02)--(2.08,.82)--(2.93,.66)--(3.78,.53)--(4.63,.43);
  \foreach \x/\y/\lab in {1.23/1.02/1,2.08/.82/2,2.93/.66/3,
                            3.78/.53/\cdots,4.63/.43/d_\ell^*} {
    \node[circle,draw=galkData,fill=white,line width=.45pt,
          minimum size=1.35mm,inner sep=0pt] at (\x,\y) {};
    \draw[galkInk,line width=.35pt] (\x,.23)--(\x,.31);
    \node[small label,anchor=north] at (\x,.20) {$\lab$};
  }
  \node[small label,anchor=west] at (5.38,.72)
    {$w_d=\alpha\rho^{d-1}$};
\end{tikzpicture}
}
  \caption{Physical feasibility and finite-horizon evaluation. (a) Coloring the conflict graph of four trajectories yields $\mathcal B_1=\{q_2,q_4\}$ and $\mathcal B_2=\{q_1,q_3\}$. Filled and open circles denote sources and destinations, respectively; stationary atom $q_0$ is excluded from the graph. (b) Candidates $A$ and $B$ have equal current costs but different resulting states. With all other physical increments equal, the illustrated relation $d_B<d_A$ gives candidate $B$ a lower next-layer transport cost. Blue and gray dashed lines denote next-layer and return trajectories, respectively; $\mathcal E$ and $\mathcal S$ denote the entanglement and storage zones. (c) Layer-distance weights decay as $w_d=\alpha\rho^{d-1}$ over the actual visible-layer set $\mathcal D_\ell^{\rm vis}$.}
  \label{fig:physical-lookahead}
\end{figure}

Future costs are estimated separately from each candidate's resulting state $s'=s_{\ell+1}$. Starting with $s_{\ell,0}=s'$, deterministic state transition $\mathcal F$ produces future states and incremental costs layer by layer:
\begin{equation}
(s_{\ell,d},\widehat C_{\ell,d})=
\mathcal F(s_{\ell,d-1},\mathcal G_{\ell+d}),
\qquad d\in\mathcal D_\ell^{\rm vis},
\label{eq:forecast-transition}
\end{equation}
where $\mathcal F$ assigns each gate in the predicted layer an unassigned legal pair of entanglement sites. It returns nonparticipating atoms that block target endpoints or trajectories to the storage zone. The unweighted physical increment at layer distance $d$ is
\begin{equation}
\widehat C_{\ell,d}=C_{\ell,d}^{\rm res}
+C_{\ell,d}^{\rm gate}+C_{\ell,d}^{\rm block},
\label{eq:forecast-components}
\end{equation}
where $C_{\ell,d}^{\rm res}$ accounts for idle-atom excitation and coherence loss of nonparticipating residents. The term $C_{\ell,d}^{\rm gate}$ accounts for physical losses from entry and temporary relocation of gate participants. The term $C_{\ell,d}^{\rm block}$ accounts for physical losses from returning blocking atoms. Both latter terms include trap-transfer, transport, and coherence losses.

The gate site for a future gate $(u,v)$ minimizes the combined travel distances of its two participants and nearest-storage distances of endpoint blockers. Ties are broken by gate-site and orientation indices. If joint entry is infeasible, the predictor tests bounded temporary relocation and then atom-by-atom entry. If neither succeeds, the candidate is marked infeasible within the prediction horizon. Its current-boundary evaluation remains available as a fallback if no candidate has a feasible prediction.

For nonempty $\mathcal D_\ell^{\rm vis}$, let $d_\ell^*=\max\mathcal D_\ell^{\rm vis}$. A terminal clearing potential $\Phi(s_{\ell,d_\ell^*})$ accounts for residency costs remaining at the end of the window:
\begin{equation}
\widehat C_\ell^{\rm future}=
\sum_{d\in\mathcal D_\ell^{\rm vis}}\alpha\rho^{d-1}\widehat C_{\ell,d}
+\alpha\rho^{d_\ell^*-1}\Phi(s_{\ell,d_\ell^*}).
\label{eq:lookahead-cost}
\end{equation}
When $H_{\max}=0$, the visible set satisfies $\mathcal D_\ell^{\rm vis}=\varnothing$ and the future cost is $\widehat C_\ell^{\rm future}=0$.

The terminal clearing potential $\Phi$ estimates the return cost remaining beyond the window. It constructs feasible returns for entanglement-zone atoms absent from the window's final layer and uses the resulting negative log fidelity increment. It is zero when the visible-layer set is empty.

\subsection{Candidate Ranking and Boundary Selection}

Future savings may require a higher current cost, so candidate selection bounds this trade-off. The solver uses a low-current-cost candidate as an anchor and limits the current-cost increase of other candidates. It then ranks them using the finite-horizon objective:
\begin{equation}
\mathbf J_\ell=(C_\ell^{\rm cur}+\widehat C_\ell^{\rm future},B_\ell,
T_\ell^{\rm rr},D_\ell^{\rm rr},\kappa_\ell)
\label{eq:lexicographic}
\end{equation}
Lexicographic ordering compares the first component, then successive components to break ties. The canonical key $\kappa_\ell=(\overline{\mathbf x}_\ell,\overline a_\ell)$ combines the chromosome and storage-site assignment to order otherwise tied candidates. Total negative log cost, time, and distance are quantized at $10^{-12}$, $10^{-6}\,\mu\mathrm{s}$, and $10^{-6}\,\mu\mathrm{m}$, respectively; batch counts use integer comparison.

Anchor construction includes gate-site projection and residency protection. After the main search, gate-site projection fixes the residency segment of evaluated candidates and adjusts gate-site genes coordinate-wise to reduce current cost. The residency encodings include the current best, all-resident, and return-or-residency recommendations obtained by comparing physical round-trip costs. Resource-limited cases scan bounded gate-site options once; other cases use at most two passes. Each trial candidate is renormalized and checked for physical feasibility. Projected candidates and previously evaluated candidates with the same residency encoding form the neighborhood $\mathcal N_\ell^{\rm guard}$. The number of additional evaluations is denoted by $E_{\rm guard}$.

Residency protection compares an atom's residency loss before visible reuse with the loss of a round trip through storage. The former includes idle-atom excitation and coherence loss; the latter includes trap-transfer, transport, and coherence losses along feasible paths. Atoms with lower residency costs that satisfy mandatory-return and capacity constraints form the protected set $\mathcal P_\ell$. Candidate selection prioritizes their residency; the set is empty when $H_{\max}=0$. For candidate $u=(\mathbf x,a)$, the number of protected atoms returned to storage is
\begin{equation}
v_\ell(u)=\sum_{q_i\in\mathcal P_\ell}
\mathbf 1\!\left[b_{\ell,i}(u)=1\right].
\label{eq:stay-violation}
\end{equation}
Let $\mathcal X_\ell^{\rm eval}\subseteq\mathcal X_\ell^{\rm phys}$ contain physically feasible candidates whose evaluation is complete. Evaluation includes the current state transition and, when look-ahead is enabled, future-layer prediction. The solver examines candidate groups in ascending $v_\ell$ and selects the first group containing a candidate with a feasible prediction. Within this group, the current physical key $\mathbf J_\ell^{\rm cur}=(C_\ell^{\rm cur},B_\ell,T_\ell^{\rm rr},D_\ell^{\rm rr},\overline{\mathbf x}_\ell)$ determines anchor $u^0$. Define $L_\ell(u)=C_\ell^{\rm cur}(u)+\widehat C_\ell^{\rm future}(u)$. Let $\Delta C_{\rm cur}(u)=\max\{0,C_\ell^{\rm cur}(u)-C_\ell^{\rm cur}(u^0)\}$ and $\Delta S_{\rm future}(u)=\max\{0,\widehat C_\ell^{\rm future}(u^0)-\widehat C_\ell^{\rm future}(u)\}$. The admission set is
\begin{equation}
\begin{aligned}
\mathcal X_\ell^{\rm adm}=\{u\in\mathcal X_\ell^{\rm eval}:{}&
v_\ell(u)=v_\ell(u^0),\\[-1mm]
&\Delta C_{\rm cur}(u)\le\eta\Delta S_{\rm future}(u),\\[-1mm]
&L_\ell(u)\le L_\ell(u^0)\}.
\end{aligned}
\label{eq:admission-domain}
\end{equation}
Here, $\eta=0.25$ limits the current-cost increase to one quarter of the predicted future savings. The admission set is a subset of the physical feasible domain.
If no candidate passes the finite-horizon feasibility check, selection uses $\mathbf J_\ell^{\rm cur}$ among current-boundary feasible candidates. Otherwise, lexicographic selection is performed over the admission set:
\begin{equation}
(\mathbf x_\ell^*,a_\ell^*)=
\mathop{\mathrm{lexarg\,min}}_{(\mathbf x,a)\in\mathcal X_\ell^{\rm adm}}
\mathbf J_\ell(\mathbf x,a).
\label{eq:boundary-objective}
\end{equation}
This rule first minimizes the sum of current and future costs. Ties are resolved by batch count, rearrangement time, transport distance, and canonical encoding, in that order.

\subsection{Exact Enumeration and Genetic Search}

Small layer-boundary problems use exhaustive enumeration, while larger problems use genetic search. The boundary encoding set is $\mathcal Z_\ell=\prod_j\{0,\ldots,|\Omega_{\ell,j}|-1\}\times\{0,1\}^{|\mathcal D_\ell|}$, with size $D_\ell=|\mathcal Z_\ell|$. All encodings are enumerated when $D_\ell\le\min(D_{\rm enum},E_{\max})$; otherwise, bounded genetic search is used~\cite{holland1975adaptation}. Here, $D_{\rm enum}$ is the enumeration threshold, and $E_{\max}$ is the evaluation budget for distinct chromosomes. Both approaches assign return storage sites using~\eqref{eq:kbest} and share the same physical evaluation and admission rules.

The deterministic seed set $\mathcal S_\ell^0$ contains all-resident, all-return, physical-greedy, and return-recommendation encodings. Chromosomes in the least-recently-used (LRU) cache with an identical complete boundary-state key receive priority. Single-point crossover is applied separately to the gate-site and residency segments. Offspring not produced by crossover undergo one of four mutations: random single-gene mutation, gate-site reselection for high-cost transport, coupled gate-site--residency mutation, or return-cost-guided residency-bit flipping. Candidates are normalized under the constraints and deduplicated by chromosome key. Cached objective values are reused for chromosomes that have already been evaluated.

Search stops at $E_{\max}$ evaluations, after $I$ generations, or after $p_{\rm stop}$ consecutive generations without improvement. If budget remains, bounded local neighborhood descent refines the current best solution. Random seeds are fixed to make the search reproducible. Let $\mathcal U_\ell=(s_\ell,\mathcal G_\ell)$ denote the current boundary instance. The parameter set $\Theta$ contains zone and AOD parameters, physical model constants, candidate-domain limits, and search budgets. Algorithm~\ref{alg:solver} gives the complete solution procedure.

\begin{algorithm}[t]
\caption{Joint Layer-Boundary Placement}
\label{alg:solver}
\scriptsize
\begin{algorithmic}[1]
\Procedure{Evaluate}{$\mathcal U_\ell,\mathbf x,\mathcal D_\ell^{\rm vis},\Theta$}
  \State $\mathcal A\gets\Call{StorageAssignment}{\mathcal U_\ell,\mathbf x}$; $\mathcal R_{\mathbf x}\gets\varnothing$
  \For{each storage-site assignment $a\in\mathcal A$}
    \State Derive blocker relocation, temporary relocation, and atom trajectories
    \If{occupancy, capacity, AOD, or intersection constraints fail}
      \State \textbf{continue}
    \EndIf
    \State Color the conflict graph; derive executable batches and resulting state $s'$
    \State Compute current cost $C^{\rm cur}$ and rearrangement metrics $B,T,D$
    \State $(\widehat C^{\rm future},\zeta)\gets\Call{FutureCost}{s',\mathcal D_\ell^{\rm vis},\Theta}$
    \State Add the current physical record and prediction-feasibility flag $\zeta$ to $\mathcal R_{\mathbf x}$
  \EndFor
  \State \textbf{return} $\mathcal R_{\mathbf x}$
\EndProcedure
\Statex
\Procedure{OptimizeBoundary}{$\mathcal U_\ell,H_{\max},\Theta$}
  \State Compute encoding bound $D_\ell$
  \State $\mathcal D_\ell^{\rm vis}\gets\Call{VisibleDepths}{\mathcal U_\ell,H_{\max},\Theta}$
  \State $\mathcal R\gets\varnothing$
  \If{$D_\ell\le\min(D_{\rm enum},E_{\max})$}
    \For{each $\mathbf x\in\mathcal Z_\ell$}
      \State $\mathcal R\gets\mathcal R\cup\Call{Evaluate}{\mathcal U_\ell,\mathbf x,\mathcal D_\ell^{\rm vis},\Theta}$
    \EndFor
  \Else
    \State $\mathcal C\gets\Call{NormalizeDedup}{\mathcal S_\ell^0,\text{cached seeds, random seeds}}$
    \For{$g=1,\ldots,I$}
      \State Generate and evaluate crossover or mutation offspring
      \State Add feasible offspring records to $\mathcal R$
      \State Retain the $P$ canonical chromosomes with the best objective values
      \If{budget is exhausted or stagnation reaches $p_{\rm stop}$ generations}
        \State \textbf{break}
      \EndIf
    \EndFor
  \EndIf
  \State Apply bounded local refinement to the current best solution in $\mathcal R$
  \If{no record in $\mathcal R$ has $\zeta=1$}
    \State \textbf{return} the record minimizing the current physical key, with $\widehat C^{\rm future}=0$
  \EndIf
  \State Merge main-search and gate-site--residency protection records
  \State Construct the admission set $\mathcal X_\ell^{\rm adm}$ from these records
  \State \textbf{return} the candidate in $\mathcal X_\ell^{\rm adm}$ minimizing $\mathbf J_\ell$
\EndProcedure
\end{algorithmic}
\end{algorithm}

\subsection{Feasibility and Solver Properties}

Candidate generation, projection, and admission selection share the constraints of $\mathcal X_\ell^{\rm phys}$. Every accepted state transition satisfies the physical constraints and ordering dependencies between rearrangement batches. Occupancy injectivity, zone capacity, AOD ordering, gate dependencies, and per-atom clock consistency are preserved across layer transitions.

If projection and admission selection are inactive and the evaluation budget covers the entire bounded candidate space, exhaustive enumeration returns the lexicographic optimum over $\mathcal X_\ell^{\rm phys}$. When these stages are active, it returns the best candidate in $\mathcal X_\ell^{\rm adm}$. Genetic search terminates within its fixed budget without a global-optimality guarantee.

\subsection{Complexity and Implementation}

The C++17 implementation precomputes coordinate distances, trajectory compatibility, and conflict edges. Deterministic caches support candidate decoding, storage-site assignment, conflict-graph coloring, and future states. Let $N_{\rm bnd}$ denote the number of boundaries between consecutive two-qubit layers. Table~\ref{tab:solver-bounds} specifies budgets by $N_{\rm bnd}$ to bound candidate evaluation, storage-site matching, and prediction depth for long circuits.

\begin{table}[!t]
\caption{Default and Size-Adaptive Solver Limits for GA-LK}
\label{tab:solver-bounds}
\centering
\PaperTableSetup
\begin{tabular*}{\columnwidth}{@{\extracolsep{\fill}}lccc@{}}
\toprule
Configuration & Enum. limit & $(E,P,I,K)$ & $H_{\rm cap}$\\
\midrule
Default & 512 & $(576,8,6,4)$ & 8\\
$N_{\rm bnd}>512$ & 64 & $(192,\le6,\le4,\le2)$ & $\le4$\\
\shortstack[l]{$N_{\rm bnd}\ge5000$\\$n\le16$} & 16 & $(32,\le4,\le1,1)$ & --\\
$N_{\rm bnd}\ge5000$ & -- & -- & $\le2$\\
\bottomrule
\end{tabular*}
\PaperTableNote{$E$, $P$, $I$, and $K$ denote the candidate evaluation budget, population size, generation count, and number of injective assignments, respectively. ``--'' imposes no additional limit in that row. Overlapping conditions apply jointly, with the tighter limit taking precedence.}
\end{table}

Here, $n=|\mathcal Q|$, and each atom has at most 6 candidate storage sites. Exact conflict-graph refinement handles at most 24 trajectories, with a branch-and-bound limit of $2\times10^5$ nodes. Linear assignment subproblems and Murty branching generate injective storage-site assignments. Constructing pairwise trajectory compatibility has quadratic cost in the number of trajectories.

Normally, $\Omega_{\ell,j}$ contains all entanglement-site pairs that pass static-occupancy and single-trajectory intersection checks. For deep circuits with few gates per layer, the first search uses a bounded gate-site domain ordered by distance. The full domain is restored if this search produces no feasible candidate. Table~\ref{tab:solver-bounds} and the code repository provide the size thresholds and complete budgets. Caches are keyed by complete boundary states and canonical encodings. Final rearrangement metrics are computed from the generated ZAIR instruction sequence.

Suppose the main search and candidate projection evaluate $E$ and $E_{\rm guard}$ candidates, respectively. Each candidate considers at most $K$ storage-site assignments and predicts $|\mathcal D_\ell^{\rm vis}|$ future layers. The approximate time complexity is
\begin{equation}
O\!\left((E+E_{\rm guard})K
\left(C_{\rm feas}+|\mathcal D_\ell^{\rm vis}|C_{\rm pred}\right)\right),
\label{eq:complexity}
\end{equation}
where $C_{\rm feas}$ and $C_{\rm pred}$ are the costs of physical feasibility evaluation and one-layer future-state estimation, respectively. Wide-layer and deep-circuit instances use bounded gate-site neighborhoods for $E_{\rm guard}$. The generated ZAIR instruction sequence must satisfy occupancy, AOD, and unintended-intersection constraints.
